\documentclass[addpoints]{exam}
% Shared exam preamble (used with \documentclass[addpoints]{exam} in each test file)
\usepackage[margin=0.5in]{geometry}
\usepackage{tikz}
\usepackage{titling}
\usepackage{calc}
\usepackage{amsmath}
\usepackage{xparse}
\usepackage{multirow}
\usepackage{listings}
\usepackage{ifthen}
\usepackage{graphicx}
\usepackage{diagbox}
\usepackage{xcolor}
\usepackage{textcomp}
\usepackage{multicol}    % <<--- needed if you use multicols in the document
\usepackage{enumitem}

\newcommand{\overbar}[1]{\overline{#1}}

% SystemVerilog listings
\lstdefinelanguage{SystemVerilog}[]{Verilog}{
  morekeywords={
    logic,byte,int,shortint,longint,bit,enum,typedef,struct,union,
    always_comb,always_ff,always_latch,unique,priority,inside,with,
    foreach,rand,randc,assert,property,sequence,disable,fork,join_none,
    package,import,export,virtual,interface,modport
  }
}
\lstset{
  language=SystemVerilog,
  basicstyle=\ttfamily\small,
  keywordstyle=\color{blue},
  commentstyle=\color{gray},
  stringstyle=\color{red},
  numbers=left,
  numberstyle=\tiny,
  stepnumber=1,
  numbersep=5pt,
  showstringspaces=false,
  tabsize=2,
  breaklines=true,
  columns=fullflexible,
  literate={//§}{{\underline{\hspace{2cm}}}}3
}


\newcounter{subp}
\setcounter{subp}{0}

% 4-variable K-map for AB\CD (no parameters to avoid # in restricted mode)
\newcommand{\kmapABCD}{%
  {\renewcommand{\arraystretch}{1.1}\small
  \begin{tabular}{|c|c|c|c|c|}\hline
  AB\textbackslash CD & 00 & 01 & 11 & 10 \\ \hline
  00 &  &  &  &  \\ \hline
  01 &  &  &  &  \\ \hline
  11 &  &  &  &  \\ \hline
  10 &  &  &  &  \\ \hline
  \end{tabular}}%
}

% 2-input truth table: headers and 4 rows for 00,01,10,11
\newcommand{\truthtable}[3]{%
  {\renewcommand{\arraystretch}{1.1}\small
  \begin{tabular}{|c|c|c|}\hline
  #1 & #2 & #3 \\ \hline
  0  & 0  &    \\ \hline
  0  & 1  &    \\ \hline
  1  & 0  &    \\ \hline
  1  & 1  &    \\ \hline
  \end{tabular}}%
}

\newcommand{\asciitable}{%
  \begingroup
  \fontsize{8}{9}\selectfont
  \centering
  \begin{tabular}{*{8}{cc}}
  00 & NUL   & 01 & SOH   & 02 & STX   & 03 & ETX   \\
  04 & EOT   & 05 & ENQ   & 06 & ACK   & 07 & BEL   \\
  08 & BS    & 09 &       & 0A & LF    & 0B & VT    \\
  0C &       & 0D & CR    & 0E &       & 0F &      \\
  10 & DLE   & 11 & DC1   & 12 & DC2   & 13 & DC3   \\
  14 & DC4   & 15 & NAK   & 16 & SYN   & 17 & ETB   \\
  18 & CAN   & 19 & EM   & 1A & SUB   & 1B & ESC   \\
  1C & FS    & 1D & GS    & 1E & RS    & 1F & US    \\
  20 & space & 21 & !    & 22 & "    & 23 & \#   \\
  24 & \$    & 25 & \%   & 26 & \&   & 27 & '    \\
  28 & (     & 29 & )    & 2A & *    & 2B & +    \\
  2C & ,     & 2D & -    & 2E & .    & 2F & /    \\
  30 & 0     & 31 & 1    & 32 & 2    & 33 & 3    \\
  34 & 4     & 35 & 5    & 36 & 6    & 37 & 7    \\
  38 & 8     & 39 & 9    & 3A & :    & 3B & ;    \\
  3C & <     & 3D & =    & 3E & >    & 3F & ?    \\
  40 & @     & 41 & A    & 42 & B    & 43 & C    \\
  44 & D     & 45 & E    & 46 & F    & 47 & G    \\
  48 & H     & 49 & I    & 4A & J    & 4B & K    \\
  4C & L     & 4D & M    & 4E & N    & 4F & O    \\
  50 & P     & 51 & Q    & 52 & R    & 53 & S    \\
  54 & T     & 55 & U    & 56 & V    & 57 & W    \\
  58 & X     & 59 & Y    & 5A & Z    & 5B & [    \\
  5C & \textbackslash & 5D & ] & 5E & \textasciicircum & 5F & \_ \\
  60 & \textasciigrave & 61 & a & 62 & b & 63 & c \\
  64 & d     & 65 & e    & 66 & f    & 67 & g    \\
  68 & h     & 69 & i    & 6A & j    & 6B & k    \\
  6C & l     & 6D & m    & 6E & n    & 6F & o    \\
  70 & p     & 71 & q    & 72 & r    & 73 & s    \\
  74 & t     & 75 & u    & 76 & v    & 77 & w    \\
  78 & x     & 79 & y    & 7A & z    & 7B & \{   \\
  7C & |     & 7D & \}   & 7E & \textasciitilde & 7F &      \\
  \end{tabular}
  \endgroup
}

\newcommand{\test}[2]{\input{../#1.tex}\newcommand{\testname}{#2}}
\newcommand{\ul}[1]{\underline{\hspace{#1cm}}}
% Test title page - includes centered title, name/NetID fields, and newpage
% Optional logo can be set with \newcommand{\courselogo}{path/to/logo.png}
\newcommand{\maketesttitle}{%
  \begin{center}
  \ifdefined\courselogo
    \includegraphics[height=1.5cm]{\courselogo} \\[1em]
  \fi
  {\Huge \textbf{\coursenum{} \\[0.5em]
  \semester{} \\[0.5em]
  \coursename{} \\[1em]
  \testname}}
  \end{center}
  \vspace{50pt}
  \noindent Name: \ul{5} \\[50pt]
  \noindent NetID: \ul{2} \\[50pt]
  \newpage
}

\graphicspath{{../img/}{../figs/}}
\newcommand{\cpp}[1]{\lstinputlisting[language=C++]{code/#1}}
\newcommand{\cuda}[1]{\lstinputlisting[language=CUDA]{code/#1}}
\newcommand{\ccode}[1]{\lstinputlisting[language=C]{code/#1}}
\newcommand{\bash}[1]{\lstinputlisting[language=bash]{code/#1}}
\newcommand{\verilog}[1]{\lstinputlisting[language=SystemVerilog]{code/#1}}
\newcommand{\python}[1]{\lstinputlisting[language=Python]{code/#1}}


\title{ECE456 Network Centric Programming -- Socket Programming}
\date{}
\author{}

\begin{document}

\maketitle

\begin{questions}

%------------------------------------------------
\question[6]

Fill in the missing line to send a short message.

\begin{lstlisting}[language=C]
#include <sys/socket.h>

int main() {

    int sock;

    const char* msg = "hello";

    ________(sock, msg, 5, 0);
}
\end{lstlisting}

\vspace{1cm}

%------------------------------------------------
\question[6]

What condition indicates an error from \texttt{recv()}?

\begin{lstlisting}[language=C]
int n = recv(sock, buf, 1000, 0);
\end{lstlisting}

Circle the correct answer.

\begin{itemize}
\item n $<$ 0
\item n $<$ 1000
\item n $>$ 0
\end{itemize}

\vspace{1cm}

%------------------------------------------------
\question[6]

What does a return value of **0** from \texttt{recv()} indicate?

Short answer.

\vspace{2cm}

%------------------------------------------------
\question[8]

What is wrong with the following code?

(single phrase)

\begin{lstlisting}[language=C]
int n = recv(sock, buf, 1024, 0);

if(n < 1024) {
    printf("error\n");
}
\end{lstlisting}

Answer:

\vspace{2cm}

%------------------------------------------------
\question[10]

Consider the following client code.

\begin{lstlisting}[language=C]
const char* msg = "HELLO";

send(sock, msg, 5, 0);
\end{lstlisting}

Is it guaranteed that all 5 bytes are sent?

\begin{itemize}
\item Yes
\item No
\end{itemize}

If no, what must the program do?

\begin{itemize}
\item retry in a loop
\item reopen the socket
\item call fork()
\end{itemize}

%------------------------------------------------
\question[10]

Fill in the missing code so the program sends the entire buffer.

\begin{lstlisting}[language=C]
int sent = 0;

while(sent < len) {

    int n = send(sock,
                 buf + sent,
                 len - sent,
                 0);

    if(n < 0)
        return -1;

    __________;
}
\end{lstlisting}

\vspace{1cm}

%------------------------------------------------
\question[10]

Fill in the missing code so the program receives a full message.

\begin{lstlisting}[language=C]
int received = 0;

while(received < len) {

    int n = recv(sock,
                 buf + received,
                 len - received,
                 0);

    if(n < 0)
        return -1;

    __________;
}
\end{lstlisting}

\vspace{1cm}

%------------------------------------------------
\question[8]

What does the following IP address represent?

\begin{center}
127.0.0.1
\end{center}

Short answer.

\vspace{2cm}

%------------------------------------------------
\question[8]

Which protocol guarantees reliable delivery?

\begin{itemize}
\item UDP
\item TCP
\end{itemize}

\vspace{1cm}

%------------------------------------------------
\question[10]

A TCP client sends a message.

The client later receives an **ACK** packet.

Does this guarantee that the server application has processed the message?

\begin{itemize}
\item Yes
\item No
\end{itemize}

\vspace{1cm}

%------------------------------------------------
\question[10]

What layer generates the ACK packet?

\begin{itemize}
\item Application
\item TCP
\item IP
\end{itemize}

%------------------------------------------------
\question[12]

Consider the following server code.

\begin{lstlisting}[language=C]
int client = accept(server_fd, NULL, NULL);

char buf[100];

recv(client, buf, 100, 0);

printf("%s\n", buf);
\end{lstlisting}

Is it guaranteed that the full message was received?

\begin{itemize}
\item Yes
\item No
\end{itemize}

Brief reason:

\vspace{2cm}

%------------------------------------------------
\question[12]

What is wrong with the following code?

(single phrase)

\begin{lstlisting}[language=C]
char buf[100];

int n = recv(sock, buf, 100, 0);

printf("%s\n", buf);
\end{lstlisting}

Answer:

\vspace{2cm}

%------------------------------------------------
\question[15]

Write a loop that reads from a socket until the connection closes.

Use:

\begin{itemize}
\item recv
\item break when recv returns 0
\end{itemize}

\vspace{6cm}

\end{questions}

\end{document}